% Section 4: Floquet eigenvalue problem of a periodicity cell, flux identity, dimension of the decaying subspaces
% (Lemma A of research/projects/bie-enclosure/theory/port-subspace.md, reviewed
% 12/12 in port-subspace-review.md), the fully discrete criterion, ordered QZ.
\section{Decaying subspaces of the half-strips}\label{sec:leads}

\subsection{The Floquet eigenvalue problem}\label{sec:pencil}
Anchor the port states of one periodicity cell at its own walls,
$z_{\mathrm L}=(a_{\mathrm L};b_{\mathrm L})$ and $z_{\mathrm R}=(a_{\mathrm R};b_{\mathrm R})$.
The scattering relation \eqref{eq:smatrix} is equivalent to
\begin{equation}\label{eq:pencil}
  \mathsf E_{\mathrm L}z_{\mathrm L}=\mathsf E_{\mathrm R}z_{\mathrm R},\qquad
  \mathsf E_{\mathrm L}=\begin{bmatrix}-R_{\mathrm L}&I\\ T_{\mathrm{LR}}&0\end{bmatrix},\qquad
  \mathsf E_{\mathrm R}=\begin{bmatrix}0&T_{\mathrm{RL}}\\ I&-R_{\mathrm R}\end{bmatrix}.
\end{equation}
In a periodic half-strip the right wall of one cell is the left wall of the
next.
A sequence of states $(z_j)_{j\ge0}$ on consecutive walls is compatible if
$\mathsf E_{\mathrm L}z_j=\mathsf E_{\mathrm R}z_{j+1}$ for all $j$. A Floquet
solution $z_{j+1}=\zeta z_j$ solves the \emph{Floquet eigenvalue problem}
\begin{equation}\label{eq:floquet}
  (\mathsf E_{\mathrm L}-\zeta\mathsf E_{\mathrm R})\,z_0=0,
\end{equation}
a generalized eigenvalue problem whose eigenvalue $\zeta$ is the Floquet
multiplier of the solution (in numerical linear algebra,
$\mathsf E_{\mathrm L}-\zeta\mathsf E_{\mathrm R}$ is the matrix pencil
$(\mathsf E_{\mathrm L},\mathsf E_{\mathrm R})$). We call the problem \emph{regular} if
$\det(\mathsf E_{\mathrm L}-\zeta\mathsf E_{\mathrm R})$ does not vanish identically in
$\zeta$. It then has $2K$ multipliers in $\CC\cup\{\infty\}$, counted with
algebraic multiplicity; $\zeta=\infty$ occurs when $\mathsf E_{\mathrm R}$ is singular. For $S_+$ we need the subspace of initial states whose
compatible sequences are square summable, and for $S_-$ the
analogous subspace towards $-\infty$, obtained from the reversed problem
$(\mathsf E_{\mathrm R}-\zeta\mathsf E_{\mathrm L})z_0=0$.

The continuous half-strip has as many decaying as growing solutions in a
gap. For the discrete problem this count is not automatic. A $2K\times2K$
generalized eigenvalue problem may have any number of multipliers inside the unit circle, and if the number is
not $K$, the closure of \cref{sec:closure} is not square. The next lemma
shows that the count follows from energy conservation of the truncated cell.

\subsection{Flux and the dimension of the decaying subspace}\label{sec:flux}
Let $u_M$ be the retained part \eqref{eq:rayleigh} of a field on a wall. Its
flux through the wall in the direction $+e_x$ is a Hermitian form in the
port state:
\begin{equation}\label{eq:flux}
  \Phi(z)=\Im\int_{I_d}\overline{u_M}\,\partial_xu_M\,\dd y=z^*\mathsf J z,
  \qquad
  \mathsf J=\begin{bmatrix}\Re\boldsymbol\gamma&-\ii\Im\boldsymbol\gamma\\
     \ii\Im\boldsymbol\gamma&-\Re\boldsymbol\gamma\end{bmatrix}.
\end{equation}
Indeed, the cross terms between different orders vanish by orthogonality. A
propagating order contributes $\gamma_m(|a_m|^2-|b_m|^2)$, and an evanescent
order contributes $2\kappa_m\Im(\bar a_mb_m)$. Each $2\times2$ block of
$\mathsf J$ has eigenvalues $\pm|\gamma_m|$. Away from Wood points
$\mathsf J$ is therefore nondegenerate with $K$ positive and $K$ negative
eigenvalues.

\begin{assumption}[truncated periodicity cell]\label{ass:cell}
$k>0$ and $\beta$ are real, every propagating order is retained, and no
retained order is at a Wood point. The scattering matrix in
\eqref{eq:pencil} is that of the exact single-cell problem: retained incident
orders, zero incident amplitude in the deleted orders, and outgoing fields in
all orders. This problem has a unique solution for every incident vector.
\end{assumption}

\begin{lemma}[dimension of the decaying subspaces]\label{lem:dimension}
Under \cref{ass:cell}, suppose the problem \eqref{eq:floquet} is regular and
$\mathsf E_{\mathrm L}-\zeta\mathsf E_{\mathrm R}$ is invertible for all $|\zeta|=1$.
Then the generalized eigenspaces of the multipliers with $|\zeta|<1$ span a
subspace $\mathcal Z_+\subset\CC^{2K}$ of dimension exactly $K$, and so do
those with $|\zeta|>1$, including $\zeta=\infty$. Both are $\mathsf J$-neutral:
$v^*\mathsf Jw=0$ for all $v,w$ in the same subspace. The square-summable
compatible sequences in $S_+$ are exactly those with $z_0\in\mathcal Z_+$.
\end{lemma}

\begin{proof}
\emph{Cell identity.} Let $h=(a_{\mathrm L};b_{\mathrm R})$ be an incident vector
and let $G_{\mathrm L}h$ and $G_{\mathrm R}h$ be the resulting states on the two
walls, with
\[
  G_{\mathrm L}=\begin{bmatrix}I&0\\ R_{\mathrm L}&T_{\mathrm{RL}}\end{bmatrix},\qquad
  G_{\mathrm R}=\begin{bmatrix}T_{\mathrm{LR}}&R_{\mathrm R}\\ 0&I\end{bmatrix}.
\]
Green's formula on the cell, with real $k^2\rho$, transmission conditions on
the rods and quasi-periodic conditions on the top and bottom edges, shows that
the flux $\Im\int\bar u\,\partial_xu\,\dd y$ of the exact field is the same on
both walls. A deleted order has a single outgoing evanescent amplitude on each
wall, so its flux is zero. Hence $\Phi(G_{\mathrm L}h)=\Phi(G_{\mathrm R}h)$ for all
$h$, and by polarization
\begin{equation}\label{eq:flux-identity}
  G_{\mathrm R}^*\mathsf JG_{\mathrm R}=G_{\mathrm L}^*\mathsf JG_{\mathrm L}.
\end{equation}
Every compatible pair $(z_{\mathrm L},z_{\mathrm R})$ is of the form
$(G_{\mathrm L}h,G_{\mathrm R}h)$, since $h$ can be read off from the first block of
$z_{\mathrm L}$ and the second block of $z_{\mathrm R}$.

\emph{Neutrality.} Let $Z$ be a basis of the generalized eigenspace of the
multipliers inside the unit circle. There is a matrix $P$ with spectral radius
less than one and $\mathsf E_{\mathrm L}Z=\mathsf E_{\mathrm R}ZP$. For $c,e\in\CC^K$ the
sequences $z_j=ZP^jc$ and $w_j=ZP^je$ are compatible, and
\eqref{eq:flux-identity} gives $z_j^*\mathsf Jw_j=z_{j+1}^*\mathsf Jw_{j+1}$. Since
$P^j\to0$, $z_0^*\mathsf Jw_0=0$. The subspace of multipliers outside the circle
is treated in the same way with the reversed recursion. Generalized
eigenvectors are included because the argument uses only $P^j\to0$.

\emph{Dimension.} In the eigenbasis of $\mathsf J$ the flux reads
$\|p\|^2-\|q\|^2$ with $p,q\in\CC^K$. On a neutral subspace, $p=0$ forces
$q=0$, so the map to $p$ is injective and the dimension is at most $K$.
Regularity and invertibility on the unit circle make the two subspaces a
direct sum equal to $\CC^{2K}$, so each has dimension $K$.

\emph{Square-summable sequences.} Project the recursion onto the subspace of
the multipliers outside the circle. Its component satisfies
$w_j=P_{\mathrm{rev}}^\ell w_{j+\ell}$ with $P_{\mathrm{rev}}^\ell\to0$, so it vanishes
for a bounded sequence.
\end{proof}

The lemma does not need the multipliers to be simple, and it allows zero and
infinite multipliers. Three remarks delimit it.
\begin{remark}[bands]\label{rem:bands}
If $k^2$ lies in a band, unimodular multipliers occur. If they are all
semisimple, the decaying subspace plus the positive-flux part of each
unimodular eigenspace again has dimension $K$. Without semisimplicity this
fails. The flux-preserving transfer matrix
$\left[\begin{smallmatrix}1+\ii&1\\1&1-\ii\end{smallmatrix}\right]$ with
$\mathsf J=\diag(1,-1)$ has a single eigenvector, of zero flux.
\end{remark}
\begin{remark}[omitted propagating orders]\label{rem:omitted}
If a propagating order is deleted, the flux of the retained orders decreases
along the half-strip by the outgoing power of that order. The count can still equal
$K$, so the count does not detect the error, but the model now loses energy.
For this reason \cref{ass:cell} requires all propagating orders to be retained;
see also \cref{sec:truncation}.
\end{remark}
\begin{remark}[Wood points]\label{rem:wood-flux}
At a retained Wood order the corresponding block of $\mathsf J$ vanishes, and
the proof breaks at its first step. In value--derivative coordinates the
threshold solution $q+p(x-x_p)$ has the nondegenerate flux $\Im(\bar qp)$.
\Cref{sec:wood} explains why this change of coordinates alone does not make
the solver regular at Wood points.
\end{remark}

\subsection{A criterion for the computed cell}\label{sec:criterion}
The computed scattering matrices are not exact, and
\eqref{eq:flux-identity} holds only approximately. The following criterion
gives the count for the computed matrices themselves. Let
$A=\mathsf E_{\mathrm L}$ and $B=\mathsf E_{\mathrm R}$ be the computed blocks and define
\begin{equation}\label{eq:defect}
\begin{gathered}
  F=G_{\mathrm R}^*\mathsf JG_{\mathrm R}-G_{\mathrm L}^*\mathsf JG_{\mathrm L},\qquad
  g=\min_{|\zeta|=1}\sigma_{\min}(A-\zeta B),\\
  j=\min_{|m|\le M}|\gamma_m|,\qquad W=G_{\mathrm R}-G_{\mathrm L}.
\end{gathered}
\end{equation}

\begin{proposition}[count for the computed cell]\label{prop:criterion}
If $0<\underline g\le g$ and $\|F\|<j\,\sigma_{\min}(W)\,\underline g$, then the
computed problem \eqref{eq:floquet} has exactly $K$ multipliers inside the unit circle and $K$
outside, including $\infty$, counted with algebraic multiplicity.
\end{proposition}

\begin{proof}[Proof sketch]
With $V_\zeta=\diag(I,\zeta I)$ and the block swap $\Pi$ one checks
$(G_{\mathrm R}-\zeta G_{\mathrm L})V_\zeta=V_\zeta\Pi(A-\zeta B)$. On the unit circle the two
matrix functions therefore have the same singular values, and in particular
$\sigma_{\min}(W)\ge g$. Set $H=(G_{\mathrm R}+G_{\mathrm L})W^{-1}$. Then
$H^*\mathsf J+\mathsf JH=2W^{-*}FW^{-1}$, so
$\Delta=\mathsf J^{-1}W^{-*}FW^{-1}$ makes $H-\Delta$ exactly $\mathsf J$-skew. Along
the homotopy $H-t\Delta$, $0\le t\le1$, the matrices move by at most
$\|F\|/(j\,\sigma_{\min}(W))<\underline g$ on the circle. They stay invertible
there, and the number of multipliers inside the circle is constant. At $t=1$
the eigenvalues of $H-\Delta$ pair as $h\leftrightarrow-\bar h$ and avoid the
imaginary axis. The Cayley map $\zeta=(h+1)/(h-1)$ sends the $K$ eigenvalues
in each half-plane to the inside and the outside of the circle. The full
proof is in \cref{app:criterion}.
\end{proof}

The separation $g$ is the distance of $A-\zeta B$ from singularity on the
circle, not the radial distance of the multipliers from the circle. The two
can differ arbitrarily for nonnormal problems, which is why the criterion uses
$g$. A computable lower bound $\underline g$ follows from $\sigma_{\min}$ at
finitely many points on the circle and a Lipschitz bound
(\cref{rem:app-gbar}). In the experiments we report the observed counts and
the flux defects, which are at rounding level; we do not evaluate the
criterion with a verified $\underline g$.

\subsection{Ordered QZ}\label{sec:qz}
An ordered QZ decomposition $\mathsf Q\mathsf E_{\mathrm L}\mathsf V=\mathsf U_{\mathrm L}$,
$\mathsf Q\mathsf E_{\mathrm R}\mathsf V=\mathsf U_{\mathrm R}$ with the $K$ multipliers
$|\zeta|<1$ first gives an orthonormal basis $\mathsf Z_+$ (the first $K$
columns of $\mathsf V$) of $\mathcal Z_+$, and
\begin{equation}\label{eq:propagation}
  \mathsf E_{\mathrm L}\mathsf Z_+=\mathsf E_{\mathrm R}\mathsf Z_+\mathsf P_+,\qquad
  \mathsf P_+=\mathsf U_{\mathrm R,11}^{-1}\mathsf U_{\mathrm L,11}.
\end{equation}
The states $z_j=\mathsf Z_+\mathsf P_+^jc$ on consecutive walls are the decaying
solutions in $S_+$. For $S_-$, the reversed problem gives $\mathsf Z_-$
and $\mathsf P_-$ with $\mathsf E_{\mathrm R}\mathsf Z_-=\mathsf E_{\mathrm L}\mathsf Z_-\mathsf P_-$. The
whole Schur subspace keeps multiple and nonnormal multipliers. A change of
basis $\mathsf Z_\pm\mapsto\mathsf Z_\pm U$ only conjugates $\mathsf P_\pm$. The solver
rejects the frequency, and reports this, if $\mathsf E_{\mathrm L}-\zeta\mathsf E_{\mathrm R}$ is
numerically singular for all $\zeta$, if the cluster size is not $K$, or if a multiplier lies in a
tolerance annulus around the circle.
